\section*{The Propers: Detailed Commentary}

\subsection*{Received mercy and the humiliation that raises\enspace\elementref{Entrance; Collect}}

Psalm 48's cultic geography is exacting: mercy is received \emph{in the midst} of God's Temple, while praise extends to earth's ends. Gathering is neither private consolation nor possession of God by one place. The antiphon lets the Church enter as recipient and witness.

Augustine's \emph{Exposition on Psalm 48} receives Zion and the temple ecclesially, placing the city's praise within the spread of the Gospel among the nations. This direct reception extends the psalm's Israelite city and worship toward the nations. Christian universality grows from a text addressed through Zion. Augustine's widening horizon presents praise to the earth's ends as mercy received at the center and carried into a world that remains God's.

The Collect's Latin axis joins the Son's \latin{humilitas}, a fallen world, release from sin's servitude, and enduring joy. Humility is therefore neither self-contempt nor mere style. It is first Christological: the Son descends and raises. Only then is it moral, because the raised people may share holy gladness.

\subsection*{A king without the instruments of domination\enspace\elementref{First Reading; Responsorial Psalm}}

Zechariah 9:9--10 holds apparently contrary predicates together. The king is just and saving, yet humble; royal, yet riding the ordinary animal of peaceful entry; victorious, yet abolishing chariot, horse, and warrior's bow. His dominion from sea to sea is not an unlimited version of coercive empire. The text interprets its scope through the peace he proclaims to the nations.

The reading's two verses form one movement. Verse 9 addresses Daughter Zion and Jerusalem with joy at the king's coming; verse 10 removes Ephraim's chariots, Jerusalem's horses, and the battle bow before announcing peace and world-spanning rule. The animal alone does not carry the political-theological reversal. The ruler is not weak because he lacks the equipment of domination. His justice and saving arrival are displayed precisely in a peace that reaches former battle centers and then the nations.

Justin Martyr's \emph{Dialogue with Trypho}, chapter 53, directly adduces Zechariah 9:9 in interpreting Jesus' entry and extends the royal sign toward the nations. His second-century apologetic argument quotes Zechariah 9:9 and reads it of Christ. Origen's \emph{Commentary on John} X.17 cites both verses and asks how the abolition of chariots and proclamation of peace follow if interpretation stops at the literal entry. He develops a spiritual account of Christ's pacifying rule. Origen comments on the verses themselves; in their literal sense, Zechariah's disarmament is public and communal.

Cyril of Jerusalem's \emph{Catechetical Lecture 12}, section 10, uses the humble king and the covenant language later in the chapter as recognizable signs of Christ. Chrysostom's \emph{Homily 66 on Matthew} contrasts the entering king with rulers surrounded by chariots and guards, drawing moral attention to his meekness. Clement of Alexandria also receives the oracle in \emph{Paedagogus} I. These Greek witnesses differ in genre---apology, commentary, catechesis, homily, moral instruction---but converge on a kingship whose recognizable form is humility. Their convergence is on Christ; what each phrase first referred to in Zechariah's own day is a separate question.

Their differences are as informative as their agreement. Justin argues from fulfillment, Origen presses the relation between the entry and the disarmament of verse 10, Cyril forms catechumens through prophetic signs, and Chrysostom draws the king's manner into moral vision. No single witness exhausts the oracle. Read cumulatively, they make it harder to detach Christological fulfillment from peace to the nations or to retain the animal while forgetting the abolished weapons. The modern pairing with Matthew 11 shares this received focus on meek kingship.

The oracle addresses Zion and gives voice to Israel's hope. Matthew's passion narrative explicitly receives the entry image elsewhere, while this Sunday pairs the oracle with Jesus' self-identification as meek. That liturgical correlation warrants Christological fulfillment while the first reading retains its original communal horizon.

Psalm 145 responds with a vocabulary of divine kingship that is ethical rather than merely titular. The Lord is gracious, compassionate, good to all, faithful, and near to those bowed down. The textual refrain names God as King; the authorized \emph{Alleluia} response instead lets paschal praise frame the same verses. Neither changes the appointed strophes. The Psalm also qualifies human hopes: the assembly praises God as King while receiving the promised human royal figure in relation to him.

\Needspace{2\baselineskip}
Augustine's \emph{Exposition on Psalm 145} emphasizes the difference between mortal rulers and the Creator who remains faithful, executes justice, feeds, frees, raises, and reigns forever. The appointed strophes do not include every one of those acts, but their full psalmic setting makes the refrain morally concrete. To call God king is to praise a rule known by compassion and fidelity. This direct psalm reception also prevents the Zecharian ruler from becoming an autonomous rival: Christian interpretation recognizes the Son's kingship within the saving action of the one God, not as a second model of domination.

\subsection*{The Spirit's dwelling and the body's future\enspace\elementref{Second Reading}}

Romans 8:9 locates Christian existence by indwelling: belonging to Christ is inseparable from the Spirit of Christ. After the appointed omission of verse 10, verse 11 grounds hope in the Spirit who raised Jesus and will give life to mortal bodies. The body is mortal and awaits vivification; it is not the evil that must be destroyed.

Verses 12--13 then announce an obligation without returning to self-salvation. Hearers are debtors, but not to disordered flesh as a ruling power. They must put evil deeds to death \emph{by the Spirit}. Chrysostom's direct exposition treats the warning as real while distinguishing sinful deeds from bodily substance. Augustine likewise keeps divine aid and human action together. Grace neither ratifies passivity nor turns obedience into autonomous achievement.

Chrysostom's \emph{Homily 13 on Romans} treats the Spirit's indwelling and the promise that the Spirit who raised Christ will vivify mortal bodies. Resurrection hope confirms embodiment rather than canceling it. \emph{Homily 14} then reads the debt and warning of verses 12--13 as a call to mortify evil actions, not to hate the body's created nature.

Augustine's \emph{Sermon 156}, sections 6--15, argues from the same command that human action is real because it is enabled: “by the Spirit” excludes both fatalistic passivity and proud self-sufficiency. Aquinas's \emph{Commentary on Romans} VIII, lectio 3, gives a medieval analysis of the debt, mortification, Spirit-led action, and the filial movement that follows just beyond the appointed verses. The latter context may explain the direction of Paul's argument, but the assembly did not hear verses 14--17 in this proclamation. These direct Latin and Greek receptions agree on the main boundary: the “flesh” to which Christians owe nothing is disordered rule and sinful deed, not the body promised resurrection.

The Lectionary omits verse 10, creating a proclamation that moves from Spirit-defined identity directly to resurrection and obligation. The complete chapter informs the reading, but worshippers do not hear the unappointed verse.

\subsection*{Revelation to the little and the yoke of the Son\enspace\elementref{Acclamation; Gospel}}

The acclamation adapts Matthew 11:25 and performs what it says: the assembly blesses the Father and receives the Gospel as gift rather than mastery. “Little ones” does not canonize ignorance or forbid disciplined study. In context it contrasts receptive dependence with the self-securing posture that refuses Jesus' works and wisdom.

Chrysostom's \emph{Homily 38 on Matthew} reads the “little ones” as simple and receptive rather than unintelligent, and he directs the saying against pride. Revelation remains the Son's gift; it is not the reward of anti-intellectualism. The acclamation's liturgical act is therefore exact: before analysis continues, the assembly praises the Father and receives the proclaimed Son. Disciplined study belongs inside that posture, not over against it.

The Gospel's center is relational before it is therapeutic. The Son receives all from the Father; only Father and Son know one another in the unique manner expressed here; those to whom the Son reveals enter by gift. The invitation to the burdened follows this claim. Rest is communion under the revealing Son, not a technique detached from his identity.

Chrysostom insists on the unique reciprocity of Father and Son and then follows the turn toward invitation: the one making the high Christological claim summons all who labor and are burdened, not a spiritually accomplished class. The same homily refuses to make the yoke effortless in the sense of having no labor; it is good and light because disciples learn meekness and humility from this teacher and are aided under his rule. The passage's “all” therefore joins universality to a definite relation with the Son.

Augustine's \emph{Sermon 19 on the New Testament} (Benedictine Sermon 69) centers the invitation on humility as the deep foundation by which human frailty is healed rather than flattered. In \emph{Confessions} VII.21 he illuminatingly reuses Matthew 11:25--29 to contrast ideas found in Platonist books with the way of the humble Christ: intellectual ascent alone did not give him the face, tears, and healing of the incarnate mediator. Augustine passes no verdict on all philosophy here, but he sharply distinguishes Christian rest from a technique of mental elevation.

Aquinas's \emph{Catena Aurea} on Matthew 11:25--30 attests only a medieval compilation practice that arranged inherited excerpts verse by verse. Its embedded fragments, not independently traced to their underlying editions and loci, remain leads and ground no patristic claim. Chrysostom's homily and Augustine's sermon carry the substantive direct reception, and the catena adds no authority beyond them.

This reception history also clarifies the Gospel's internal order. Thanksgiving in verses 25--26 is addressed to the Father; the Son's unique knowledge and power to reveal are stated in verse 27; only then does the universal invitation sound in verses 28--30. Humility is consequently not a method for bypassing Christology, and rest is not offered by an anonymous wisdom teacher. The meek heart belongs to the Son who reveals the Father. Conversely, the high claim of mutual knowledge does not terminate in exclusion: it opens into “come to me” for all who labor. Chrysostom and Augustine keep the doctrinal claim and the ascetical invitation together, neither swallowing the other.

The same order controls “learn from me.” Jesus does not ask disciples to admire meekness from a distance but to enter a teacher--learner relation embodied in a yoke. Learning can include study, correction, practice, and perseverance; the promise of rest describes the goodness of belonging to this master, not exemption from formation. Chrysostom's warning against pride and Augustine's foundation image converge here as two distinct interpretations. Receptivity makes serious learning possible because knowledge of the Father is received through the Son rather than seized as a possession.

The yoke image prevents an interpretation of rest as absence of claim. A yoke joins a learner to directed labor; Jesus calls hearers to take and learn. Its goodness lies in the teacher's meek and humble heart and in the fittingness of his lordship. Christian meekness therefore cannot mean acquiescing in abuse or refusing duties of protection. It is strength governed by truth and charity under Christ.

Ambrose's \emph{Exposition of the Christian Faith} V.167 receives the yoke as a glorious subjection borne in humility and gentleness. Leo the Great, in \emph{Letter 14}, applies the meek and humble Master to ecclesial leadership, warning pastors against placing unbearable burdens on those they govern. These are illuminating doctrinal and pastoral reuses, not verse-by-verse Gospel commentaries. Their applications expose a test internal to the text: an authority claiming Christ's yoke while multiplying crushing burdens contradicts the teacher whose heart defines that yoke.

\begin{studybox}{The Fathers and Doctors on each text}
Augustine expounds Psalms 48, 145, and 34; Justin, Origen, Cyril, Clement, and Chrysostom cite Zechariah's king and read it of Christ; Chrysostom, Augustine, and Aquinas expound Romans 8; and Chrysostom and Augustine expound Matthew 11. Ambrose, Leo, and Augustine's \emph{Confessions} apply the yoke and the invitation in doctrine, pastoral rule, and autobiography rather than in verse-by-verse comment. Aquinas's \emph{Catena} attests only medieval compilation practice, and its embedded fragments, not independently traced, remain leads.
\end{studybox}

\subsection*{Purification and heavenly conduct\enspace\elementref{Prayer over Offerings}}

The Prayer over the Offerings moves from a dedicated oblation to purification and heavenly conduct. It does not identify the gifts with moral effort or make purification the congregation's purchase. God receives the Church's offering within Christ's action and is asked to transform those who offer. Because the prayer is common to A/B/C, relating it to the meek yoke is a synthesis from its place beside this Gospel, not a Lectionary correlation.

\subsection*{Two Communion paths and one final petition\enspace\elementref{Communion A; Communion B; Prayer after Communion}}

Communion Antiphon A appoints Psalm 34's invitation to taste and see within sacramental reception. Its full psalmic context joins tasting to blessing, seeking, deliverance, instruction, and refuge. It does not reduce Eucharistic truth to a pleasant feeling; objective gift and fitting disposition remain distinct.

Augustine's \emph{Exposition on Psalm 34} treats tasting as knowledge gained through participation and presses the invitation toward life and conduct rather than curiosity at a distance. He reads the psalm of Christ and the Church centuries after its composition, and that reading explains why the verse became durable Eucharistic language: the good is not merely discussed but received. The full psalm's fear of the Lord, guarded speech, turning from evil, and pursuit of peace supplies the moral horizon; sacramental tasting cannot be isolated from conversion.

\Needspace{3\baselineskip}
Alternative B repeats Matthew 11:28 and gives this cycle a particularly transparent Gospel--Communion correspondence. On this path, the burdened hear the same invitation at proclamation and sacramental approach. The two appointed antiphons remain distinct alternatives, each with its own textual contribution.

The Prayer after Communion asks that recipients filled with great gifts receive salvation and continue in praise. Its order matters: divine gift precedes enduring response. “Unceasing” praise cannot mean continuous speech; it names a life ordered toward God, including work, rest, service, repentance, and the next Eucharistic gathering.

\section*{The Propers: Notable and Quotable}

{\small
\begin{itemize}
\item \textbf{“My yoke is easy and my burden light” --- a spectacularly indecent burden.\enspace\elementref{Gospel}} In James Gillray's 1782 satire of the Worsley scandal, Sir Richard physically supports his wife's lover so the man can look through a bathhouse window. He announces Christ's line while bearing the voyeur on his shoulders. The visual pun makes “yoke” bodily, and the sacred promise becomes savage ridicule of complaisance and humiliation.

\item \textbf{“My burden is light” --- a novel titled for a burden that is anything but.\enspace\elementref{Gospel}} Agatha Christie's official site confirms that her 1956 Mary Westmacott novel \emph{The Burden} prints Matthew 11:30 as its epigraph. Its fierce sisterly story blurs love, jealousy, obsession, family, and identity, letting the promised lightness sit tensely over a psychologically crushing form of love.

\item \textbf{“My yoke is easy and my burden is light” --- a sweetheart supplies the weight.\enspace\elementref{Gospel}} An 1862 California newspaper prints the line as a young man's joke while his sweetheart sits on his lap with her arm around his neck. “Yoke” becomes the embrace and “burden” the woman herself: a sacred paradox is compressed into a flirtatious physical pun.
\end{itemize}
}

\section*{The Propers: Interpretive Possibilities}

\noindent\emph{The following exploratory proposals were developed in an AI-assisted editorial process. They are attributed to none of the cited authorities and claim no historical compositional intent. Each states its anchors and limits.}

\Needspace{6\baselineskip}
\noindent\textbf{The King replaces the war machine with a school.\enspace\elementref{First Reading; Gospel}}
Zechariah removes chariot, warhorse, and bow; Matthew gives yoke, teacher, and learning. The mechanism changes authority's instrument from force to pedagogy---a conversion missed separately. Its fruit tests political, ecclesial, parental, and intellectual rule by whether it forms truthful agency. Justin's colt bearing Christ's word-yoke is a genuine ancient analogue. Yet Matthew 11 does not cite Zechariah: neither authorial design nor the illegitimacy of every coercive restraint follows.

\Needspace{6\baselineskip}
\noindent\textbf{Disarmament reaches the habits that can rebuild the bow.\enspace\elementref{First Reading; Second Reading}}
Zechariah removes weapons; semi-continuous Romans mortifies bodily deeds by the indwelling Spirit. The mechanism makes institutional and interior disarmament correct each other: separate reading misses how systems return through desire, fear, domination, and vengeance. Its fruit joins justice to conversion without substitution. Paul is not commenting on Zechariah, “flesh” is not the body, and privatizing war, excusing aggressors, or blaming victims disconfirms the analogy.

\Needspace{6\baselineskip}
\noindent\textbf{Revelation to the little becomes resurrection from within.\enspace\elementref{Second Reading; Acclamation; Gospel}}
The acclamation and Gospel give the little ones revelation of the Father; Romans gives the indwelling Spirit who vivifies bodies. The mechanism is participatory knowing: the gift inhabits and remakes the knower. Separate reading leaves an anti-intellectual slogan beside an isolated datum; the fruit is rigorous, humble study tested in prayer and embodied change. Hilary and Irenaeus closely join indwelling to renewal, but no checked source joins this Pauline account to Matthew's little ones. It fails against expertise or by confusing the Son's revelation with the Spirit's indwelling.

\Needspace{9\baselineskip}
\noindent\textbf{The one raised from collapse is entrusted with a fitted weight.\enspace\elementref{Responsorial Psalm; Gospel; Communion B}}
Psalm 145 raises the bowed; the Gospel gives rest, then Christ's yoke; Communion B repeats the summons only if selected. The mechanism moves from rescue through restored capacity to freely borne, shared discipline. Separate reading misses rehabilitated agency; its fruit resists earned worth and permanent infantilization. The branch remains unresolved; invoking the yoke to preserve abuse, deny medical or legal relief, rush trauma recovery, or disguise a manipulator's burden as divine command rules the proposal out.

\Needspace{6\baselineskip}
\noindent\textbf{Two Communion branches propose two ways of learning the good.\enspace\elementref{Gospel; Communion A; Communion B; Prayer after Communion}}
The Gospel says “learn”; A, if chosen, answers by tasting and seeing; B, if chosen, repeats “come”; the final prayer tests knowledge in continuing praise. Compared as alternatives, the two branches reveal participation and reiterated call---branch architecture missed separately. Cyprian's yoke--tasting juxtaposition is a near analogue, not this liturgical argument. The fruit is learning verified in durable life. Neither branch makes feeling grace's measure, guarantees worthy reception, or licenses both antiphons for one undocumented celebration.
